\paragraph{Identité de la section et composition de ses observables.}
\label{el-ampl-articulacion}
L'identité structurelle de l'électron est introduite par la section
centrale et ses opérations de transport. L'inversion de la coordinatisation APP
sélectionne la cellule $(5,5)$ par sa propriété de point fixe. Les deux feuillets
arithmétiques permettent d'étudier les déformations qui conservent la somme et
le reste du produit. Le résultat distingue les deux orientations
$(8,2)$ et $(2,8)$ et conserve la différence d'une unité de quotient.
Tel est le premier niveau d'information : une cellule, une involution, une
fibre conservée et un défaut entier minimal. Aucune évaluation de masse
n'est nécessaire pour l'énoncer ou le démontrer.

Le passage aux projections arithmétiques introduit un deuxième niveau. Les
fibres de somme et de produit déterminent des espérances conditionnelles sur le
même espace fini. Leur commutateur mesure la dépendance à l'égard de
l'ordre de ces deux projections. La norme $\sqrt3/4$ résulte du spectre
de leur appariement ; le bloc qui atteint cette norme conserve en outre
les relations quadratiques démontrées ci-dessous. Le
préfacteur scalaire et l'algèbre locale ont donc une provenance commune, mais
contiennent des informations différentes. La norme retient la grandeur du
commutateur, tandis que les matrices conservent son orientation et sa loi de
composition.

Le transfert surfacique--volumique incorpore les coordonnées régionales
à un troisième niveau. Le calendrier des modes détermine la matrice
d'incidence et son action quadratique. La marque de clôture produit le rapport
$\pi/\varphi^2$ ; l'échange des deux arguments fournit l'
orientation $\varphi/\pi^2$ utilisée dans le caractère exponentiel
électronique. Les deux expressions sont reliées par une involution
explicite. Cette relation explique pourquoi une coordonnée régionale peut
apparaître dans l'équation électronique avec une puissance ou dans une position
différente de celle qu'elle occupait dans la première lecture : l'opération appliquée
fait partie de la donnée mathématique conservée.

Le terme cubique d'alpha appartient au complémentaire régional d'un retour
marqué. Le dénombrement de ses positions produit $27-5=22$, tandis que les
trois feuillets de chacune des cinq régions produisent $3\cdot5=15$.
Les démonstrations conservent le domaine de ces dénombrements et l'
orientation qui détermine leurs signes. La formule basale compose
donc une norme d'incompatibilité, un caractère de
transfert, un déficit de positions et une mémoire torsionnelle.
L'écriture sur une seule ligne exprime la composition de ces
opérations ; son explication exige de les reconstruire dans cet ordre.

La transition de la coordonnée basale à la coordonnée complète conserve
la même section centrale. Le facteur de retour utilise les registres
de l'histoire, et les deux opérateurs de lecture déclarés restituent le triplet
$(80,54,6)$ sur cette trajectoire. La multiplication par
$R_{\mathrm{act}}^{80-54\alpha+6\Delta_4}$ modifie l'évaluation,
mais maintient identifiés son origine et ses antécédents. La valeur
basale constitue une étape intermédiaire de la composition ; la valeur
complète inclut également cette contribution du retour. La comparaison
physique doit utiliser l'étape spécifiée par la formule évaluée.

L'échelle d'action intervient ensuite avec une autre fonction. Le rapport de
deux sections d'action est sans dimension et annule leur échelle commune.
La réalisation absolue utilise, en revanche, une base d'action, un
opérateur de lecture temporel et l'application vers l'énergie ou la masse. Cette distinction permet
de conserver simultanément la construction de l'échelle et l'identité
de la section électronique. Le quantum chronologique fournit une
référence d'échelle ; la section centrale détermine le facteur structurel
exprimé par rapport à elle. Les équations dimensionnelles de ce
chapitre maintiennent cette composition explicite.

La pertinence de l'électron dans l'argument général provient de cette
convergence d'opérations. La même arithmétique qui a produit les
régions détermine maintenant une cellule centrale, ses déformations et ses
projections. La constante de structure fine, construite auparavant, agit
dans cette composition comme coordonnée de clôture. L'information de
retour détermine l'évaluation complète, et la réalisation dimensionnelle
permet la comparaison ultérieure. Expliquer l'électron à partir de la structure
consiste à conserver cette succession d'objets, d'applications et de valeurs à chaque
étape de la démonstration.
